\section{Black Holes and Hawking Radiation}
\begin{quote}\small
\noindent\textbf{Status.} \emph{Legacy physical-regime block inside the active main body.}\par
\noindent\textbf{Block function.} This section reads collapse, horizon behavior, and Hawking-type radiation as frame- and readout-level consequences of the same closure architecture.
\end{quote}
\label{sec:bh}
%% ============================================================

\begin{definition}[Schwarzschild horizon in Sphaera coordinates]
The Schwarzschild horizon $r_S=2GM/c^2$ corresponds to $\beta\to 1$ in
the Sphaera metric $g_{tt}=1-\beta^2$.
In sector coordinates: the horizon is where $\zminus\to\infty$.
\end{definition}

\begin{proposition}[Horizon as cascade singularity]
\label{prop:horizon-cascade-singularity}
The cascade operator $\KC$ is undefined at $\zminus=\infty$: the Sphaera
cannot project beyond its own horizon.
The horizon is the boundary of the Sphaera's jurisdiction.
\begin{proof}
$\KC=C^\mu{}_\nu\circ\Pi_J$ requires $\|\zminus\|<\infty$
(finite sector displacement).
At $\zminus\to\infty$: $\beta\to 1$, $\gamma_{\mathrm{Lorentz}}\to\infty$,
the Lorentz transformation is singular, and the projection $\Pi_J$ ceases
to be bounded.
\end{proof}
\end{proposition}

\begin{theorem}[Hawking radiation from the triolic structure]
\label{thm:hawking-triolic}
At the horizon, a virtual matter--tachyon pair from the triolic structure
(Theorem~3.2) is promoted to a real pair:
\begin{itemize}
  \item Tachyon ($\beta=-\bsym$): falls into the black hole.
  \item Matter ($\beta=+\bsym$): escapes to infinity.
  \item Their compensator output $C(f)=\tfrac{1}{2}(f+Jf)$: a photon
        with $\beta=0$, $\zminus=0$ --- the Hawking radiation.
\end{itemize}
The temperature of the Hawking photon in the photon frame is:
\[
  T_{\mathrm{Sphaera}}
  = \gstring\cdot\frac{m_P^2 c^2}{8\pi M k_B}
  = \gstring\cdot T_{\mathrm{Hawking}}.
\]
\end{theorem}

\begin{proof}
The Hawking temperature (standard derivation) is
\[
T_H = \frac{\hbar c^3}{8\pi G M k_B} = \frac{m_P^2 c^2}{8\pi M k_B}.
\]
The Sphaera produces the Hawking photon in the photon frame where
$\gamma$ is measured as $\gstring$ (not $\gBI$, which is the
matter-frame value).
Hence $T_{\mathrm{Sphaera}} = \gstring\cdot T_H$.
The ratio $T_{\mathrm{Sphaera}}/T_{\mathrm{Hawking}}=\gstring
=\ln 2/(\pi\sqrt{3})\approx 0.1274$ is constant across all black hole
masses.
\end{proof}

\begin{proposition}[Hawking temperature frame analysis and physical meaning of $\gBI$]
\label{prop:hawking-gBI-tachyon}
The Hawking temperature $T_{\mathrm{Sphaera}} = \gstring\cdot T_H^{(\mathrm{matter})}$
has the following frame interpretations:
\begin{enumerate}[label=(\roman*)]
  \item \textbf{Photon frame:} $k(T) = k(\hbar) + k(\kappa) - k(k_B)$
    where surface gravity $\kappa = c^4/(4GM)$ has
    $k(\kappa) = 3k(c) - k(G_N) = -2$ and $k(\hbar) = 1$,
    giving $k(T_H) = -1$ (assuming $k(k_B)=0$).
    Hence $T_H^{(\mathrm{photon})} = \gL \cdot T_H^{(\mathrm{matter})}$,
    and the Sphaera prediction is
    \[
      T_{\mathrm{Sphaera}} = \frac{\gstring}{\gL}\cdot T_H^{(\mathrm{photon})}
      = \frac{\gstring^2}{\gBI}\cdot T_H^{(\mathrm{photon})}.
    \]
  \item \textbf{Tachyon frame:} The tachyon sector has opposite frame shift,
    $T_H^{(\mathrm{tachyon})} = \gL^{-1}\cdot T_H^{(\mathrm{matter})}$. Then:
    \[
      T_{\mathrm{Sphaera}}
      = \gstring\cdot T_H^{(\mathrm{matter})}
      = \gstring\cdot\gL\cdot T_H^{(\mathrm{tachyon})}
      = \gBI\cdot T_H^{(\mathrm{tachyon})}.
    \]
  \item \textbf{Physical meaning of $\gBI$:} The Barbero--Immirzi parameter
    is the ratio of the Sphaera Hawking temperature to the tachyon-frame
    Hawking temperature:
    \[
      \gBI = \frac{T_{\mathrm{Sphaera}}}{T_H^{(\mathrm{tachyon})}}.
    \]
    This gives $\gBI$ a direct kinematic interpretation: it measures how much
    the tachyon sector amplifies the Hawking temperature relative to the
    Sphaera output. Since the tachyon falls into the black hole and the
    Sphaera photon escapes, $\gBI$ encodes the frame mismatch between the
    infalling and escaping components of the triolic virtual pair.
\end{enumerate}
\end{proposition}

\begin{proof}
Statement~(i): with $k(\hbar)=1$, $k(G_N)=2$, $k(c)=0$, $k(k_B)=0$,
the frame correction to $T_H = \hbar c^3/(8\pi G_N M k_B)$ is
$\gL^{k(\hbar)-k(G_N)} = \gL^{-1}$, giving
$T_H^{(\mathrm{photon})} = \gL\cdot T_H^{(\mathrm{matter})}$.

Statement~(ii): the tachyon sector sits at $\zminus = -\bsym$, the mirror
image of the matter frame. The frame factor for the tachyon is
$\gL^{(\mathrm{tachyon})} = 1/\gL^{(\mathrm{matter})}$ by the symmetric
matter--tachyon frame theorem~\ref{thm:symmetric-frame}. Hence
$T_H^{(\mathrm{tachyon})} = \gL^{-1}\cdot T_H^{(\mathrm{matter})}$, and
$T_{\mathrm{Sphaera}}/T_H^{(\mathrm{tachyon})}
= \gstring\cdot\gL = \gBI$.

Statement~(iii) follows immediately from (ii).
\end{proof}

\begin{quote}\small
\noindent\textbf{Reader cue.} Read the next block as a geometric recoding of the same Hawking data. Its purpose is to shift from the horizon/frame-temperature dictionary to the coupling interpretation before the section turns to the core/horizon dual picture.
\end{quote}

\begin{theorem}[Geometric origin of $\gL$ from Hawking geometry]
\label{thm:dual-origin}
\emph{(Corrected v2.29.1: previous version used incorrect $\gL\approx 1.197$
and an erroneous $2/3$ sector factor.)}
The splitting factor $\gL=\gBI/\gstring\approx 1.864$ satisfies:
\[
  \gL \;\approx\; \frac{1}{\bsym} \;=\; \frac{1}{\sin(\theta_H/2)},
\]
where $\theta_H\approx 66.6^\circ$ is the Hopf angle of the matter observer
on $S^3$. Numerically: $1/\bsym \approx 1.820$ vs.\ measured
$\gL\approx 1.864$, a discrepancy of $2.4\%$.
\end{theorem}

\begin{proof}
The matter observer at $\bsym$ has radial position $r/r_S=1/(1-\bsym^2)$
relative to the Schwarzschild horizon.
The local Hawking temperature is $T(r)=T_H(1-r_S/r)^{-1/2}=T_H/\bsym$.
Since $\gBI$ is determined by requiring $S=A/(4\lP^2)$ for the observed
temperature, the effective coupling shifts:
$\gBI = \gstring\cdot(T_H/T(r))^{-1} = \gstring/\bsym$.
Hence $\gL = \gBI/\gstring = 1/\bsym \approx 1.820$.

The residual $2.4\%$ gap between $1/\bsym=1.820$ and the measured
$\gL=1.864$ is the same continuum-limit correction as in the $\bsym$-Hopf
gap (Remark~\ref{rem:bsym-hopf-gap}): the discrete prime lattice
$\Gamma_P$ shifts $\bsym$ from the Hopf candidate $1/\sqrt{3}$ by
the factor $\delta_{\mathrm{CL}}$, and the same correction applies here.
\end{proof}

\begin{remark}[Connection to GOE/GUE spectral ratio]
\label{rem:dual-origin-goe-gue}
By Proposition~\ref{prop:goe-fermionic} and
Remark~\ref{rem:goe-gue-string-lqg},
$\gL = \gBI/\gstring$ is the ratio of the GOE scale ($H^+_-$, LQG) to
the GUE scale ($H^{++}$, String). The geometric relation $\gL\approx
1/\bsym = 1/\sin(\theta_H/2)$ therefore says:
\[
  \frac{\text{GOE scale}}{\text{GUE scale}}
  \;\approx\;
  \frac{1}{\sin(\text{matter Hopf angle}/2)}.
\]
The matter observer's position on $S^3$ (the Hopf angle $\theta_H$)
encodes the ratio of the two random-matrix universality classes.
\end{remark}

\begin{quote}\small
\noindent\textbf{Reader cue.} Read the next block as an interpretive deepening of the same closure architecture. Its purpose is to shift from the frame- and coupling-readout picture to the core/horizon communication picture without introducing a separate physical mechanism.
\end{quote}

\begin{theorem}[Hawking radiation as the voice of the core]\label{thm:core}
The Hawking tachyon is not created at the horizon.
It is the black hole core --- which has $\zminus\to\infty$ and cannot
be projected by $\KC$ --- attempting to reach $\zminus=0$ through the
only available channel: tunnelling through the horizon.

The two descriptions are CPT-dual and physically identical:
\begin{center}
\renewcommand{\arraystretch}{1.2}
\begin{tabular}{ll}
\toprule
Outside view ($t$ forward) & Inside view ($t$ reversed)\\
\midrule
$t_1$: vacuum & $t_3$: core at $\zminus=\infty$\\
$t_2$: virtual pair at horizon & $t_2$: tachyon at horizon\\
$t_3$: tachyon falls in, matter escapes & $t_1$: Hawking photon emitted\\
\bottomrule
\end{tabular}
\end{center}
The cascade $\KC$ has no preferred time direction: it is the equilibrium
between $\zminus=0$ (outside) and $\zminus=\infty$ (core).
The Hawking pair is the discharge of this tension.

\begin{proof}[Structural argument]
Outside: $\zminus$ grows from $0$ (far away) to $\infty$ (horizon).
$\KC$ attempts to project $\zminus\to 0$ but fails at the horizon.
Tension forces a real pair.
Inside: $\zminus=\infty$ (core, singular compression).
$\KC$ attempts $\zminus\to 0$ but the horizon blocks outward projection.
The only escape: a core fraction tunnels as a tachyon
($\beta=-\bsym$) through the horizon.
Under $J$-reflection, this tachyon appears outside as matter
($\beta=+\bsym$).
The compensator $C(f)=\tfrac{1}{2}(f+Jf)$ produces the Hawking photon
at $\beta=0$.

CPT\@: $C$ exchanges matter and tachyon; $P$ exchanges $\beta\leftrightarrow-\beta$
(sides of the horizon); $T$ reverses creation and annihilation.
Both descriptions are related by $CPT$.
\end{proof}
\end{theorem}

\begin{remark}[Black hole as inverted Sphaera]
The normal Sphaera has $\zminus$ decreasing inward:
$\KC$ projects from large $\zminus$ toward $\zminus=0$
(the photon ground state).
A black hole inverts this: $\zminus$ increases inward, from
$\zminus=0$ (far away) to $\zminus=\infty$ (horizon and core).
The cascade $\KC$ cannot function inside the horizon; the interior is
a region where the Sphaera does not exist.
The horizon is where the Sphaera is maximally stressed.

Hawking radiation is the Sphaera discharging at its boundary.
The black hole evaporates because the core is the Sphaera --- and the
Sphaera always tends toward $\zminus=0$.
The tachyon is not a particle falling in. It is the core, speaking
through the only open channel.
\end{remark}

\begin{remark}[The entropy factor]
The Bekenstein--Hawking entropy $S_{BH}=A/(4\lP^2)$ and the LQG entropy
$S_{LQG}=A\ln 2/(8\pi\gBI\lP^2)$ agree when $\gBI=\ln 2/(2\pi)\approx 0.110$.
The measured value $\gBI\approx 0.2375$ differs by a factor $\approx 2.15$.
Within the present framework: the LQG formula is derived in the matter
frame, while the Bekenstein--Hawking formula is a photon-frame statement.
The discrepancy factor $2.15$ should be related to the frame
transformation, but its precise form --- unlike the analogous Hawking
temperature factor $\gstring$ --- is not yet derived.
This is an open quantitative problem.
\end{remark}

\begin{theorem}[ART--QFT consistency of the Sphaera]\label{thm:artqft}
The Sphaera biquaternion encodes both General Relativity and Quantum
Field Theory, connected by Wick rotation:
\[
  it\cdot\mathbf{1} + z_i e_i
  \xrightarrow{t\mapsto i\tau_E}
  \tau_E\cdot\mathbf{1} + z_i e_i.
\]
At the natural observer position $r_0/r_S = 1/(1-\bsym^2)$, the
Schwarzschild metric gives $g_{tt}(r_0)=-\bsym^2$ (invariant): no
additional geometric factor is hidden in the ART side.
The ART--QFT conversion is internally consistent.
\end{theorem}

%% ============================================================

